\chapter{MARCO TEÓRICO}
\label{cap:marco}

En este capítulo se presentan los fundamentos conceptuales que sustentan el
desarrollo del proyecto. Se revisan el área de conocimiento, las tecnologías de
desarrollo móvil empleadas, los conceptos de arquitectura de software, los
servicios en la nube, los fundamentos de inteligencia artificial y realidad
aumentada aplicados al producto, las bases fisiológicas de la prescripción del
ejercicio y la metodología de desarrollo adoptada.

% =====================================================================
\section{Área de conocimiento}
\label{sec:mt-area}

\subsection{Programación móvil}

\subsubsection*{Definición}

La programación móvil es la rama de la ingeniería de software dedicada al
diseño, construcción y mantenimiento de aplicaciones que se ejecutan en
dispositivos portátiles, principalmente teléfonos inteligentes y tabletas.
A diferencia del software de escritorio, una aplicación móvil debe operar con
recursos limitados de batería, memoria y conectividad, y aprovechar sensores
propios del dispositivo, como la cámara, el GPS o el acelerómetro
\citep{sommerville2011}.

\subsubsection*{Características}

\begin{itemize}
    \item \textbf{Movilidad y contexto:} el usuario utiliza la aplicación en
    cualquier lugar, en este caso, dentro del gimnasio y entre series.
    \item \textbf{Recursos restringidos:} la interfaz y la lógica deben ser
    eficientes en consumo de memoria, red y batería.
    \item \textbf{Conectividad variable:} la aplicación debe tolerar
    desconexiones temporales y sincronizar datos cuando la red vuelva.
    \item \textbf{Interacción táctil:} la interfaz se diseña para pantallas
    pequeñas y gestos, con elementos de tamaño adecuado para el dedo.
    \item \textbf{Distribución mediante tiendas:} las aplicaciones se publican
    en tiendas como Google Play, que imponen requisitos de empaquetado y de
    permisos.
\end{itemize}

\subsubsection*{Elementos: tipos de aplicaciones móviles}

Según la forma de construirlas, las aplicaciones móviles se clasifican en tres
grandes grupos, resumidos en la Tabla~\ref{tab:tipos-apps}.

\begin{table}[H]
\centering
\caption{Tipos de aplicaciones móviles según su construcción}
\label{tab:tipos-apps}
\begin{tabularx}{\textwidth}{|L{2.8cm}|Y|Y|}
\hline
\rowcolor{encabezado} \textbf{Tipo} & \textbf{Descripción} & \textbf{Ventajas y desventajas} \\
\hline
Nativa & Se desarrolla con el lenguaje y SDK de cada plataforma (Kotlin para
Android, Swift para iOS). & Máximo rendimiento y acceso al hardware; requiere
una base de código por plataforma. \\ \hline
Híbrida & Aplicación web empaquetada dentro de un contenedor nativo
(\emph{WebView}). & Rápida de construir; rendimiento y experiencia inferiores a
una nativa. \\ \hline
Multiplata\-forma nativa & Un solo código fuente que se traduce a componentes
nativos de cada plataforma (React Native, Flutter). & Una sola base de código
con interfaz nativa y buen rendimiento; depende del ecosistema del
\emph{framework}. \\
\hline
\end{tabularx}
\fuente{Elaboración propia con base en \citet{reactnative} y \citet{sommerville2011}.}
\end{table}

\subsection{Sistemas de información}

Un sistema de información es un conjunto de componentes interrelacionados que
recolectan, procesan, almacenan y distribuyen información para apoyar la toma
de decisiones y el control \citep{laudon2016}. Sus actividades básicas son la
\emph{entrada} de datos, su \emph{procesamiento} para convertirlos en
información útil y la \emph{salida} hacia las personas que la necesitan,
complementadas con una \emph{retroalimentación} que permite evaluar y corregir
la entrada.

En el presente proyecto, la aplicación es un sistema de información personal:
recibe los datos del usuario (objetivo, nivel, peso, estatura), los procesa
(generación de rutinas, cálculo del IMC, agregación del historial) y devuelve
información accionable (la rutina del día, la categoría de IMC, la tendencia
del peso), cerrando el ciclo cuando el usuario registra nuevas mediciones.

% =====================================================================
\section{Desarrollo móvil multiplataforma}
\label{sec:mt-multiplataforma}

\subsection{React Native}

React Native es un \emph{framework} de código abierto creado por Meta que
permite construir aplicaciones móviles nativas con JavaScript o TypeScript y el
modelo de componentes de React \citep{reactnative}. A diferencia de las
aplicaciones híbridas, los componentes declarados (\archivo{View},
\archivo{Text}, \archivo{Image}) se traducen en vistas nativas de Android o
iOS. Su arquitectura actual, basada en la interfaz JSI y el renderizador
\emph{Fabric}, comunica el código JavaScript con el nativo de forma síncrona y
sin serialización intermedia.

\subsection{Expo}

Expo es una plataforma construida sobre React Native que provee un conjunto de
módulos nativos mantenidos (cámara, imágenes, fuentes, enlaces profundos,
pantalla de bienvenida), herramientas de línea de comandos y servicios de
compilación en la nube (EAS) \citep{expo}. Cada versión del SDK de Expo fija un
conjunto compatible de bibliotecas, lo que reduce los conflictos de versiones.
Cuando la aplicación necesita módulos nativos adicionales se utiliza un
\emph{development build}, una versión de desarrollo propia de la aplicación
que reemplaza a la aplicación genérica Expo Go.

\subsection{Expo Router}

Expo Router es el sistema de navegación oficial de Expo, basado en el sistema
de archivos: cada archivo dentro de la carpeta \archivo{app/} se convierte en
una ruta de la aplicación \citep{exporouter}. Los grupos de rutas, escritos
entre paréntesis (por ejemplo \archivo{(auth)}), organizan las pantallas sin
afectar la URL, y los archivos \archivo{\_layout.tsx} definen la navegación
contenedora (pila, pestañas). Desde el SDK~53 incorpora las \emph{rutas
protegidas} (\archivo{Stack.Protected}), que muestran u ocultan grupos de
pantallas según una condición booleana, como la existencia de una sesión.

\subsection{TypeScript}

TypeScript es un superconjunto de JavaScript desarrollado por Microsoft que
añade tipado estático opcional \citep{typescript}. El compilador verifica los
tipos antes de la ejecución y detecta errores como propiedades inexistentes o
argumentos incorrectos, lo que resulta especialmente útil cuando los datos
provienen de servicios externos con estructuras complejas.

\subsection{Bibliotecas complementarias}

\begin{itemize}
    \item \textbf{NativeWind:} adapta las clases utilitarias de Tailwind CSS a
    los estilos de React Native, de modo que la interfaz se estiliza con
    clases como \archivo{bg-primary} o \archivo{text-foreground} a partir de
    una paleta centralizada \citep{nativewind}.
    \item \textbf{Zustand:} biblioteca minimalista de gestión de estado global
    basada en \emph{hooks}, con soporte de persistencia
    \citep{zustand}.
    \item \textbf{Zod:} biblioteca de declaración y validación de esquemas que
    infiere los tipos de TypeScript a partir del esquema, de modo que la misma
    definición valida los datos en tiempo de ejecución y los tipa en tiempo de
    compilación \citep{zod}.
\end{itemize}

% =====================================================================
\section{Arquitectura de software}
\label{sec:mt-arquitectura}

\subsection{Arquitectura cliente--servidor}

En el modelo cliente--servidor, los procesos se dividen en clientes que
solicitan servicios y servidores que los proveen \citep{tanenbaum2007}. En una
aplicación móvil, el cliente es la aplicación instalada en el teléfono, que se
encarga de la presentación y de parte de la lógica, mientras que el servidor
centraliza la autenticación, la persistencia de los datos y la distribución de
contenidos multimedia. Esta separación permite que los datos de un usuario
estén disponibles en cualquier dispositivo y que el contenido compartido, como
un catálogo, se actualice sin publicar una nueva versión de la aplicación.

\subsection{Patrón Modelo--Vista--Controlador}

El patrón Modelo--Vista--Controlador (MVC), descrito originalmente para
Smalltalk-80 por \citet{krasner1988}, separa una aplicación interactiva en tres
responsabilidades:

\begin{itemize}
    \item \textbf{Modelo:} representa los datos y las reglas de acceso a ellos.
    \item \textbf{Vista:} presenta la información al usuario y captura sus
    acciones.
    \item \textbf{Controlador:} recibe las acciones de la vista, aplica la
    lógica de negocio y coordina al modelo.
\end{itemize}

La separación reduce el acoplamiento: la interfaz puede rediseñarse sin tocar
las reglas de negocio, y la fuente de datos puede cambiar sin afectar a la
interfaz. \citet{martin2017} generaliza esta idea en la \emph{regla de
dependencia}: las capas externas (interfaz, bases de datos, servicios) dependen
de las internas (reglas de negocio), y nunca a la inversa.

\subsection{Patrón repositorio}

El patrón repositorio media entre la lógica de negocio y la capa de acceso a
datos mediante una interfaz similar a una colección de objetos del dominio
\citep{fowler2002}. El resto del sistema pide ``las rutinas del usuario'' sin
conocer en qué colección, tabla o servicio se almacenan. En este proyecto, los
repositorios son el único lugar que conoce la estructura real de la base de
datos.

% =====================================================================
\section{Plataforma Firebase}
\label{sec:mt-firebase}

Firebase es una plataforma de Google que ofrece servicios administrados para
aplicaciones móviles y web bajo el modelo de \emph{backend como servicio}
(BaaS): el desarrollador consume autenticación, base de datos y almacenamiento
a través de SDK cliente, sin desplegar servidores propios \citep{firebase}.

\subsection{Firebase Authentication}

Servicio de identidad que gestiona el registro, el inicio de sesión y la
recuperación de contraseñas con distintos proveedores (correo y contraseña,
Google, entre otros). Cada usuario recibe un identificador único (\emph{uid})
que se adjunta a sus solicitudes y que las reglas de seguridad utilizan para
autorizar el acceso a los datos.

\subsection{Cloud Firestore}

Cloud Firestore es una base de datos NoSQL orientada a documentos. Los datos se
organizan en \emph{colecciones} que contienen \emph{documentos}, y cada
documento es un conjunto de pares campo--valor que puede contener, a su vez,
\emph{subcolecciones} \citep{firestore}. Sus características relevantes para
el proyecto son:

\begin{itemize}
    \item \textbf{Escuchas en tiempo real:} mediante \archivo{onSnapshot}, el
    cliente recibe automáticamente cada cambio de los documentos que observa,
    sin volver a consultar.
    \item \textbf{Persistencia sin conexión:} el SDK mantiene una caché local
    que permite leer y escribir sin red y sincroniza al reconectarse.
    \item \textbf{Marcas de tiempo del servidor:} el valor
    \archivo{serverTimestamp()} se reemplaza por la hora del servidor al
    confirmar la escritura, evitando depender del reloj del teléfono.
    \item \textbf{Subcolecciones por propietario:} los datos que pertenecen a
    un único usuario pueden anidarse bajo su documento
    (\archivo{usuarios/\{uid\}/...}), de modo que la propia ruta expresa la
    pertenencia.
\end{itemize}

\subsection{Cloud Storage para Firebase}

Servicio de almacenamiento de objetos para archivos de gran tamaño, como
imágenes y animaciones. Los archivos se identifican por una ruta dentro del
\emph{bucket} y se obtienen mediante una URL de descarga generada a partir de
esa ruta.

\subsection{Reglas de seguridad}

Las credenciales de configuración de Firebase viajan dentro de la aplicación
instalada y son, por diseño, públicas. La protección de los datos recae en las
\emph{reglas de seguridad}, un lenguaje declarativo que el servidor evalúa en
cada lectura y escritura \citep{firestorerules}. Las reglas pueden comprobar la
identidad del solicitante (\archivo{request.auth.uid}), la ruta del documento y
el contenido de los datos que se intentan escribir
(\archivo{request.resource.data}), lo que permite validar tipos y rangos en el
servidor además de en el cliente.

% =====================================================================
\section{Inteligencia artificial}
\label{sec:mt-ia}

\subsection{Definición y enfoques}

\citet{russell2020} definen la inteligencia artificial (IA) como el estudio de
los agentes que perciben su entorno y actúan para alcanzar sus objetivos de la
mejor manera posible. Dentro de la disciplina coexisten dos grandes enfoques:

\begin{itemize}
    \item \textbf{IA simbólica o basada en conocimiento:} representa
    explícitamente el conocimiento de un dominio mediante reglas, hechos y
    estructuras lógicas, y lo aplica mediante un mecanismo de inferencia. Su
    razonamiento es transparente y explicable.
    \item \textbf{Aprendizaje automático:} el sistema infiere un modelo a partir
    de ejemplos, sin que las reglas se programen explícitamente. Incluye el
    aprendizaje profundo con redes neuronales \citep{goodfellow2016}.
\end{itemize}

\subsection{Sistemas basados en reglas}

Un sistema basado en reglas, o sistema experto, es un programa que emula el
razonamiento de un especialista humano en un dominio acotado
\citep{giarratano2005}. Sus componentes son:

\begin{itemize}
    \item \textbf{Base de conocimiento:} conjunto de reglas de producción del
    tipo \emph{SI condición ENTONCES acción}, obtenidas de expertos o de la
    literatura del dominio.
    \item \textbf{Memoria de trabajo:} los hechos del caso concreto, por ejemplo,
    el objetivo y el nivel de un usuario.
    \item \textbf{Motor de inferencia:} selecciona y aplica las reglas cuyas
    condiciones se cumplen sobre los hechos para obtener una conclusión.
\end{itemize}

Este enfoque es adecuado cuando el conocimiento del dominio está
consolidado y documentado, como ocurre con las pautas de prescripción del
ejercicio publicadas por el American College of Sports Medicine, y cuando se
requiere que cada recomendación pueda justificarse. Además, no necesita un
conjunto de datos de entrenamiento, del que el proyecto no dispone.

\subsection{Visión por computadora en dispositivos móviles}

La clasificación de imágenes con redes neuronales convolucionales (CNN) asigna
a una fotografía una etiqueta de un conjunto predefinido con una probabilidad
asociada \citep{goodfellow2016}. Arquitecturas ligeras como MobileNet
\citep{howard2017} reducen el costo computacional para ejecutarse en el
teléfono. TensorFlow Lite, actualmente llamado LiteRT, es el entorno de
ejecución de Google para modelos de aprendizaje automático en dispositivos
móviles \citep{tflite}; permite cargar un modelo entrenado previamente y
ejecutar la inferencia localmente, sin enviar la imagen a un servidor.
Entrenar un clasificador de este tipo requiere un conjunto de imágenes
etiquetadas con permiso de uso para cada clase a reconocer.

% =====================================================================
\section{Realidad aumentada}
\label{sec:mt-ra}

\citet{azuma1997} define la realidad aumentada (RA) como un sistema que
combina lo real y lo virtual, es interactivo en tiempo real y registra los
objetos virtuales en el espacio real. \citet{milgram1994} la ubican dentro de
un continuo que va del entorno completamente real al completamente virtual,
donde la RA corresponde a la vista del mundo real enriquecida con elementos
generados por computadora.

Según la técnica de registro, la RA móvil se clasifica en:

\begin{itemize}
    \item \textbf{Basada en marcadores:} reconoce una imagen o código impreso y
    dibuja el contenido sobre él.
    \item \textbf{Sin marcadores (con seguimiento del entorno):} detecta
    superficies y puntos característicos para anclar objetos en tres
    dimensiones; en Android se implementa con ARCore \citep{arcore}.
    \item \textbf{Basada en la cámara con superposición de información:}
    muestra la imagen en vivo de la cámara y superpone sobre ella información
    contextual del objeto enfocado, sin anclaje tridimensional. Es la forma
    más ligera de RA y no requiere hardware certificado.
\end{itemize}

En el proyecto se implementa la tercera modalidad: el escáner muestra la
cámara en vivo, un marco de encuadre y, una vez identificada la máquina, su
ficha superpuesta sobre la imagen real.

% =====================================================================
\section{Fundamentos del entrenamiento de fuerza}
\label{sec:mt-entrenamiento}

\subsection{Principios del entrenamiento}

La prescripción del ejercicio se rige por principios ampliamente aceptados
\citep{acsm2021, kraemer2004}:

\begin{itemize}
    \item \textbf{Especificidad:} las adaptaciones dependen del tipo de
    estímulo; entrenar para fuerza, hipertrofia o resistencia exige variables
    distintas.
    \item \textbf{Sobrecarga progresiva:} para seguir progresando, el estímulo
    debe aumentar gradualmente.
    \item \textbf{Individualización:} la carga debe ajustarse a la experiencia,
    la condición física y los objetivos de cada persona.
    \item \textbf{Recuperación:} los grupos musculares necesitan entre 48 y 72
    horas de descanso entre sesiones intensas \citep{garber2011}.
\end{itemize}

\subsection{Variables de prescripción}

Las principales variables de una rutina de fuerza son la frecuencia semanal, la
selección de ejercicios, el número de series, el número de repeticiones y el
tiempo de descanso entre series. La Tabla~\ref{tab:variables-objetivo} resume
los rangos recomendados según el objetivo.

\begin{table}[H]
\centering
\caption{Variables de prescripción según el objetivo del entrenamiento}
\label{tab:variables-objetivo}
\begin{tabularx}{\textwidth}{|L{3.3cm}|C{1.6cm}|C{2.8cm}|C{2.4cm}|Y|}
\hline
\rowcolor{encabezado} \textbf{Objetivo} & \textbf{Series} & \textbf{Repeticiones} &
\textbf{Descanso} & \textbf{Observación} \\
\hline
Hipertrofia (ganar masa) & 3--6 & 6--12 & 60--120 s & Cargas moderadas a
altas \\ \hline
Resistencia muscular & 2--3 & 15--25 & 30--60 s & Cargas bajas, descansos
cortos \\ \hline
Salud / mantenimiento & 2--4 & 8--15 & 60 s & Mínimo dos sesiones por semana \\ \hline
Pérdida de grasa & 2--4 & 12--20 & 30--60 s & Combinar con trabajo
cardiovascular \\
\hline
\end{tabularx}
\fuente{Elaboración propia con base en \citet{kraemer2004},
\citet{schoenfeld2010} y \citet{acsm2021}.}
\end{table}

\subsection{División del entrenamiento}

La forma en que se reparten los grupos musculares a lo largo de la semana se
denomina división (\emph{split}). Las más comunes son el entrenamiento de
cuerpo completo, adecuado para principiantes con pocas sesiones semanales, y
la división empuje--tirón--pierna (\emph{push--pull--legs}), que agrupa los
músculos que empujan (pecho, hombros, tríceps), los que tiran (espalda,
bíceps) y el tren inferior, permitiendo más volumen por grupo con la
recuperación adecuada \citep{kraemer2004}. La frecuencia recomendada aumenta
con la experiencia: dos a tres sesiones semanales para principiantes, tres a
cuatro para intermedios y cuatro a seis para avanzados.

\subsection{Técnica y prevención de lesiones}

La ejecución correcta de cada ejercicio (postura, rango de movimiento,
velocidad controlada, ajuste de la máquina a la anatomía del usuario) es
condición para que el entrenamiento sea eficaz y seguro. Las máquinas de
resistencia guiada reducen la exigencia de estabilización y son recomendables
para iniciar, siempre que se ajusten correctamente el asiento, los topes y los
apoyos \citep{acsm2021}.

\subsection{Índice de masa corporal}

El índice de masa corporal (IMC) es un indicador de la relación entre el peso y
la estatura, utilizado como aproximación poblacional al estado nutricional
\citep{oms2000}. Se calcula como:

\begin{equation}
\mathrm{IMC} = \frac{\text{peso (kg)}}{\text{estatura (m)}^2}
\label{eq:imc}
\end{equation}

La Tabla~\ref{tab:imc-oms} presenta la clasificación de la Organización
Mundial de la Salud. El IMC no distingue entre masa muscular y grasa, por lo
que en personas muy musculadas debe interpretarse con cautela.

\begin{table}[H]
\centering
\caption{Clasificación del IMC en adultos según la OMS}
\label{tab:imc-oms}
\begin{tabular}{|L{5cm}|C{4cm}|}
\hline
\rowcolor{encabezado} \textbf{Clasificación} & \textbf{IMC (kg/m$^2$)} \\
\hline
Bajo peso & $< 18.5$ \\ \hline
Peso normal & $18.5 - 24.9$ \\ \hline
Sobrepeso & $25.0 - 29.9$ \\ \hline
Obesidad & $\geq 30.0$ \\
\hline
\end{tabular}
\fuente{\citet{oms2000}.}
\end{table}

% =====================================================================
\section{Metodología Mobile-D}
\label{sec:mt-mobiled}

Mobile-D es una metodología ágil propuesta por el centro VTT de Finlandia para
el desarrollo de aplicaciones móviles \citep{abrahamsson2004}. Combina
prácticas de Programación Extrema (programación en ciclos cortos, integración
continua, refactorización), de Crystal (escalabilidad para equipos pequeños) y
del Proceso Unificado (ciclo de vida por fases). Está pensada para equipos de
menos de diez personas y proyectos de pocos meses. Su ciclo se divide en las
cinco fases de la Figura~\ref{fig:mobiled}.

\begin{figure}[H]
\centering
\begin{tikzpicture}[node distance=0.35cm,
  fase/.style={rectangle, rounded corners=4pt, draw=black!70, fill=#1,
  minimum width=2.45cm, minimum height=1.25cm, align=center, font=\footnotesize\bfseries}]
  \node[fase=green!10] (e) {Exploración};
  \node[fase=green!18, right=of e] (i) {Inicialización};
  \node[fase=green!26, right=of i] (p) {Producción};
  \node[fase=green!34, right=of p] (s) {Estabilización};
  \node[fase=green!42, right=of s] (t) {Pruebas del\\sistema};
  \foreach \a/\b in {e/i,i/p,p/s,s/t} \draw[flecha] (\a) -- (\b);
  \draw[flechad] (p.south) to[out=-60,in=-120,looseness=4] node[below,font=\scriptsize]{iteraciones} (p.south);
  \node[below=1.1cm of e, text width=2.5cm, align=center, font=\scriptsize]
    {Alcance, requerimientos, plan};
  \node[below=1.1cm of i, text width=2.5cm, align=center, font=\scriptsize]
    {Entorno, arquitectura, día de prueba};
  \node[below=1.1cm of s, text width=2.5cm, align=center, font=\scriptsize]
    {Integración y refactorización};
  \node[below=1.1cm of t, text width=2.5cm, align=center, font=\scriptsize]
    {Verificación y corrección};
\end{tikzpicture}
\caption{Fases de la metodología Mobile-D}
\label{fig:mobiled}
\fuente{Adaptado de \citet{abrahamsson2004}.}
\end{figure}

\begin{itemize}
    \item \textbf{Exploración:} se identifican los interesados, se define el
    alcance, se recogen los requerimientos y se planifica el proyecto.
    \item \textbf{Inicialización:} se prepara el entorno técnico, se define la
    arquitectura y se realiza un \emph{día de prueba} para validar que las
    herramientas funcionan de extremo a extremo.
    \item \textbf{Producción:} se implementa la funcionalidad en iteraciones;
    cada una se compone de un \emph{día de planificación}, \emph{días de
    trabajo} y un \emph{día de liberación} en el que se verifica lo construido
    con pruebas de aceptación.
    \item \textbf{Estabilización:} se integran los módulos, se refactoriza y se
    completa la documentación, sin agregar funcionalidad nueva.
    \item \textbf{Pruebas del sistema:} se verifica que el producto cumpla los
    requerimientos y se corrigen los defectos encontrados.
\end{itemize}

% =====================================================================
\section{Calidad del software}
\label{sec:mt-calidad}

La norma ISO/IEC 25010 define un modelo de calidad del producto de software
con ocho características: adecuación funcional, eficiencia de desempeño,
compatibilidad, usabilidad, fiabilidad, seguridad, mantenibilidad y
portabilidad \citep{iso25010}. En este proyecto el modelo se utiliza como
referencia para formular los requerimientos no funcionales y para organizar
las pruebas del sistema.

% =====================================================================
\section{Herramientas de desarrollo}
\label{sec:mt-herramientas}

\begin{table}[H]
\centering
\caption{Herramientas utilizadas en el desarrollo}
\label{tab:herramientas}
\begin{tabularx}{\textwidth}{|L{3.6cm}|Y|}
\hline
\rowcolor{encabezado} \textbf{Herramienta} & \textbf{Uso en el proyecto} \\
\hline
Visual Studio Code & Editor de código principal, con extensiones de
TypeScript, ESLint, Prettier y Tailwind. \\ \hline
Android Studio & SDK de Android, emulador y compilación del
\emph{development build}. \\ \hline
Node.js y npm & Entorno de ejecución de las herramientas y gestor de
dependencias. \\ \hline
Git y GitHub & Control de versiones con ramas \archivo{main},
\archivo{develop} y \archivo{feature/*}, y solicitudes de integración
(\emph{pull requests}). \\ \hline
Consola de Firebase & Configuración de Authentication, Firestore, Storage y
publicación de reglas de seguridad. \\ \hline
ESLint y Prettier & Análisis estático y formato uniforme del código. \\ \hline
Expo CLI y EAS & Servidor de desarrollo, compilación y generación de los
paquetes de instalación. \\
\hline
\end{tabularx}
\end{table}
